\documentclass[conference]{IEEEtran}
\IEEEoverridecommandlockouts

\usepackage{cite}
\usepackage{amsmath,amssymb,amsfonts}
\usepackage{graphicx}
\usepackage{xcolor}
\usepackage{url}
\usepackage{balance}

% Figures live in repo results/figures/
\graphicspath{{../results/figures/}}

\begin{document}

\title{Informed Bit-Position Pruning for Streaming Hyperdimensional Computing on Zynq\\
{\footnotesize A Verified Silicon Study of Accuracy, Energy, and Cross-Subject Mask Transfer}
}

\author{
\IEEEauthorblockN{Harshavardhan Reddy Narra}
\IEEEauthorblockA{Ming Hsieh Department of Electrical and Computer Engineering\\
University of Southern California\\
Los Angeles, CA, USA\\
hnarra@usc.edu}
\and
\IEEEauthorblockN{Srinivas Boppu}
\IEEEauthorblockA{School of Electrical Sciences\\
IIT Bhubaneswar\\
Bhubaneswar, India\\
\textit{email TBD}}
}

\maketitle

% IEEE: avoid symbols and math in the abstract.
\begin{abstract}
Hyperdimensional computing (HDC) is attractive for low-latency edge inference, yet
prior FPGA studies often vary hypervector dimension rather than asking which bit
positions matter under compression.
We present a bit-exact streaming 1024-bit HDC classifier on Xilinx Zynq-7020,
controlled over AXI4-Lite and fed by AXI4-Stream with DMA batching, and validate it
on EMG hand-gesture recognition under a frozen five-subject protocol.
Over 658 thousand test windows, silicon replay matches the Python golden reference
with zero accuracy delta.
We map accuracy, out-of-context area, and measured board energy across dimension,
bundle precision, and Fisher-informed pruning, and show that informed pruning to
87.5 percent sparsity is accuracy-neutral on silicon while random masks at the same
density lose 8.6 percentage points.
A pooled Fisher mask trained on three subjects transfers to two held-out subjects with
only 0.86 percentage points loss.
Programmable-logic batch inference is about 200 times faster and 175 times lower in
measured energy per window than ARM software HDC on the same board.
\end{abstract}

\begin{IEEEkeywords}
Hyperdimensional computing, FPGA, Zynq, pruning, Fisher mask, EMG classification, energy measurement
\end{IEEEkeywords}

\section{Introduction}
\label{sec:intro}

% TODO: 3 paragraphs + bullet contributions
% P1: HDC on FPGA for edge EMG; prior work prunes D not bit positions.
% P2: Gap: need verified deployment + measured Pareto + iso-density pruning study.
% P3: What we built on Zynq (streaming, bit-exact, INA219 energy).

Hyperdimensional computing (HDC) offers lightweight inference for wearable sensing,
but most FPGA-HDC work reports accuracy at fixed dimension rather than studying
\emph{which} hypervector bits can be removed without harming classification.
We target EMG hand-gesture recognition on a Zynq-7020 ZedBoard with a streaming
AXI-DMA datapath and a discriminability-guided pruning mask programmed from software.

\textbf{Contributions.}
\begin{itemize}
\item A bit-exact 1024-bit HDC accelerator with XOR bind, cyclic permute, majority
bundle, and masked Hamming associative memory, verified by nine co-simulation harnesses
and 658{,}004-window EMG replay on silicon (74.24\% spatial accuracy, $\Delta=0.00\%$).
\item A Hook~A accuracy--area--energy study over dimension, bundle precision, and
Fisher keep ratio, with measured INA219 energy at silicon anchors A/B/C and an ARM
software baseline ($\sim$12 vs.\ $\sim$2088\,$\mu$J/window batch-amortized).
\item Twist~1: at iso-density, Fisher-informed masks preserve accuracy while random
masks lose 8.63 percentage points at 128 kept bits (keep ratio 0.125).
\item Twist~2: a pooled Fisher mask from three training subjects transfers to two
held-out subjects with 0.86 percentage points gap versus a local oracle mask.
\end{itemize}

\section{Background and Related Work}
\label{sec:related}

% TODO: HDC primitives; Rahimi FPGA-HDC EMG; pruning prior art (dimension vs bits).
% Cite: Rahimi ICRC 2016, recent FPGA-HDC surveys.

\subsection{Hyperdimensional Computing}
Binary spatter codes represent symbols as high-dimensional binary vectors; binding
(XOR), permutation, and bundling (majority vote) compose structured hypervectors for
classification by Hamming nearest-prototype search.

\subsection{FPGA-HDC and Pruning}
Prior FPGA-HDC EMG work established competitive accuracy with hardware acceleration.
Our work differs by (i) a verified streaming AXI path, (ii) bit-\emph{position} pruning
via Fisher masks at fixed density, and (iii) measured Zynq board energy rather than
proxy metrics alone.

\section{System Architecture}
\label{sec:arch}

% TODO: block diagram figure; cite rtl/ modules.
% Key facts: 100 MHz PL, zero DSP/BRAM, pruning_mask.sv, item_mem LUT ROM.

Figure~\ref{fig:arch} shows the datapath.
The processing system programs prototypes and the global pruning mask over AXI4-Lite;
inference windows stream over AXI4-Stream with scatter-gather DMA batching.

\begin{figure}[t]
\centering
% TODO: replace with architecture diagram (from slides or new TikZ)
\fbox{\parbox{0.92\columnwidth}{\centering\vspace{1.2cm}
\textbf{[TODO: System block diagram]}\\
PS (AXI-Lite) + PL core + DMA stream + pruning mask\vspace{1.2cm}}}
\caption{Streaming HDC system on Zynq-7020: PS programs masks and prototypes; PL
executes encode--classify; DMA delivers batched EMG windows.}
\label{fig:arch}
\end{figure}

Post-route utilization on xc7z020 for the DMA path is 35{,}206 LUTs (66.2\%) and
27{,}639 FFs, with zero DSP and zero BRAM blocks.

\section{Verification and Methodology}
\label{sec:method}

\subsection{Protocol and Encoders}
We use frozen protocol \textbf{P-may2026}: five UCI EMG configuration subjects,
stratified train/test split, seed~1.
Python Stage~B reproduces literature-class $\sim$90\% accuracy with a four-channel
record encoder; the deployed RTL encoder (Eq.~3.1 grid, seed-42 item memory) targets
74.24\% and is the silicon verification path---not a failure to match literature on
hardware.

\subsection{Co-Simulation and Board Replay}
Each RTL block is checked bit-for-bit against a Python golden model.
Board PASS requires $|$accuracy$_{\mathrm{board}} - $accuracy$_{\mathrm{ref}}| \le 0.5\%$.

\subsection{Energy Measurement}
Whole-board 12\,V input power is sensed at the ZedBoard J21 shunt with an INA219 on
a Raspberry Pi (I$^2$C).
Batch integration scales energy by measured batch duration ($\sim$0.93\,ms for PL,
$\sim$164\,ms for ARM over 200 windows), not the full logging interval.

\subsection{Hook A, Twist 1, and Twist 2}
Hook~A sweeps $D \in \{256,512,1024,2048\}$, CNT\_W $\in \{3,4,5,6\}$, keep
$\in \{1.0,0.5,0.25,0.125\}$ with per-subject Fisher masks (320 rows).
Silicon anchors A/B/C use a pooled Fisher mask at $D{=}1024$, CNT\_W${=}6$.
Twist~1 compares informed vs.\ random masks at equal kept-bit count (five random seeds).
Twist~2 trains a pooled mask on subjects 1--3 and evaluates on held-out subjects 4--5.

\section{Results}
\label{sec:results}

\subsection{Verification and Baselines}
Table~\ref{tab:baselines} summarizes deployment baselines at matched protocol.
PL DMA batch inference is $\sim$200$\times$ faster than ARM software HDC; measured
energy is $\sim$175$\times$ lower batch-amortized.

\begin{table}[t]
\caption{Deployment baselines (P-may2026, five subjects).}
\label{tab:baselines}
\centering
\begin{tabular}{lccc}
\hline
\textbf{Path} & \textbf{Acc.} & \textbf{Latency} & \textbf{Energy} \\
\hline
PL DMA batch & 74.24\% & $\sim$4\,$\mu$s/w & 11.98\,$\mu$J/w \\
ARM HDC & 74.15\% & 819\,$\mu$s/w & 2088\,$\mu$J/w \\
Tiny int8 MLP & 93.0\% & --- & --- \\
\hline
\end{tabular}
\end{table}

\subsection{Hook A: Accuracy--Area--Energy Pareto}
At $D{=}1024$ and CNT\_W $\ge 4$, informed Fisher pruning is flat at 74.15\% from
0\% to 87.5\% prune in Python; silicon anchor replays A/B/C remain $\sim$74.3\%.
Measured total energy at anchors A/B/C is similarly flat ($\sim$12\,$\mu$J/w),
indicating static-dominated board power rather than prune-dependent joules.

\begin{figure}[t]
\centering
\includegraphics[width=\linewidth]{hookA_pareto_measured.png}
\caption{Hook~A Pareto with measured silicon energy: (a) Python OOC LUT sweep with
placed anchor; (b) measured $\mu$J/window at anchors A/B/C and ARM path.}
\label{fig:pareto}
\end{figure}

\begin{figure}[t]
\centering
\includegraphics[width=\linewidth]{fisher_heatmap.png}
\caption{Fisher score heatmap and informed mask cutoffs at keep ratios 0.5 and 0.125.}
\label{fig:fisher}
\end{figure}

\subsection{Twist 1: Informed vs.\ Random at Iso-Density}
At keep${=}0.125$ (128 bits), informed masks retain 74.15\% spatial mean accuracy
while random masks average 65.51\% ($+8.63$\,pp gap, $\ge$5\,pp target met).
At keep${=}0.5$ the gap is smaller ($+1.70$\,pp), consistent with random masks
retaining enough signal at moderate density.

\begin{figure}[t]
\centering
\includegraphics[width=\linewidth]{twist1_informed_vs_random_keep0125.png}
\caption{Twist~1 headline at keep${=}0.125$: informed vs.\ random iso-density pruning.}
\label{fig:twist1}
\end{figure}

\subsection{Twist 2: Cross-Subject Mask Transfer}
A pooled Fisher mask from subjects 1--3 yields 68.45\% on held-out subjects 4--5
versus 69.31\% for local oracle masks ($+0.86$\,pp gap, below 3\,pp pilot threshold).

\begin{figure}[t]
\centering
\includegraphics[width=\linewidth]{twist2_cross_subject.png}
\caption{Twist~2: local oracle vs.\ pooled transfer mask on held-out subjects.}
\label{fig:twist2}
\end{figure}

\section{Discussion and Limitations}
\label{sec:limit}

Measured board energy does not decrease with pruning because total $\mu$J/window is
static-dominated at batch throughput; pruning primarily reduces effective search width
and area proxy, not J21 input joules under our measurement setup.
The deployed RTL encoder ($\sim$74\%) intentionally differs from literature Stage~B
($\sim$90\%) encoding; claims are relative to the verified deployment path.
Twist experiments use five configuration subjects; broader leave-one-subject-out
studies are future work.

\section{Conclusion}
\label{sec:conclusion}

We presented a verified streaming HDC accelerator on Zynq with measured accuracy,
area, and energy characterization.
Informed Fisher bit-position pruning preserves accuracy through aggressive compression
where random iso-density masks fail, and pooled masks transfer across subjects with
small loss.
Future work includes PL-rail power isolation, larger subject cohorts, and encoder
studies to narrow the deployment vs.\ literature accuracy gap.

\section*{Acknowledgment}
% TODO: funding, lab, dataset credit (HDC-EMG GPLv3 fetch).

\bibliographystyle{IEEEtran}
\bibliography{refs}

\end{document}
